\begin{definecredits}
	\hrspercred{30}
	\wkspercred{15}
	\Lcredmul{1}
	\Tcredmul{2}
	\Pcredmul{2}
% LO	-3:0:0
% LOb	-2:0:0
% LOb	-1:0:0
% LT	-3:1:0
% LP	-3:0:1
% PO	-0:0:1
%NPTEL	-2:0:0
% Proj	-0:0:5  --Minor Project
% ProjMaj	-0:0:15  --Major Project
% ProInt	-0:0:5  --Industry Internship/Research Internship/ Projects in CoEs
%%%%%%%%%%%%%%%%%%%%%%%%%%%%%%%%%%%%%%%%%%%%%%%%%%
%-----------NPTEL Courses----------------
	\begin{credit}{NPTEL}
	\Lcredit{2} 
	\SEEduration{3}
	\SEELmarks{100}
	\guideline{Students shall register for the corresponding NPTEL course on the SWAYAM portal and complete all weekly assignments within the stipulated deadlines.}
	\guideline{Continuous evaluation is based on the NPTEL platform's own weekly assignment scores, and the semester-end certification exam conducted by NPTEL.}
	\guideline{A minimum aggregate score, as prescribed by NPTEL for the elective/certificate course, must be secured to be eligible for course credit.}
	\guideline{Attendance and completion records will be monitored by the concerned Course Coordinator.}
	\end{credit}
	\begin{credit}{NPTEL}
		\variantchar{RM}
		\Lcredit{2}  
		\SEEduration{3}
		\SEELmarks{100}
		\guideline{NPTEL is an acronym for National Programme on Technology Enhanced Learning, which is an initiative by seven Indian Institutes of Technology (IIT Bombay, Delhi, Guwahati, Kanpur, Kharagpur, Madras, and Roorkee) and the Indian Institute of Science (IISc) for creating course contents in engineering and science.}
		
		\guideline{NPTEL offers online certification courses through its portal:
			\url{https://swayam.gov.in/nc_details/NPTEL}}
		
		\guideline{Enrollment for courses and examination registration can be done in ONLINE mode only. The link is available on the NPTEL website:
			\url{http://nptel.ac.in/}}
		
		\guideline{Students need to enroll for the NPTEL course and successfully clear the examination.}
		
		\guideline{In case a student fails to obtain the certificate, the student needs to enroll for the same course again in the subsequent NPTEL semester and clear the examination.}
		
		\guideline{If the same course is not offered by NPTEL (i.e., if the same course is not re-run) in the subsequent semester, the student needs to write a letter seeking permission from the Counsellor, HoD, and Dean Academics, with further approval from the BoS Committee, to take an alternative course from the list announced by NPTEL.}
		
		\guideline{The examination shall be conducted by NPTEL.}
	\end{credit}


%-----------Minor Project Courses----------------
\begin{credit}{Proj}
	\Lcredit{0}
	\Pcredit{5}
	\CIEduration{2} 
	\SEEduration{3} 
	\CIELmarks{100}
	\SEELmarks{100}
	\commoninstCIE{The evaluation committee shall consist of Guide, Professor, Associate Professor/Assistant
		Professor. The committee shall assess and evaluate the presentation and the progress reports.}
	\commoninstSEE{The SEE examination shall be conducted by an external examiner (domain expert) and an internal
		examiner. \textbf{Evaluation shall be carried out in batches, not exceeding 6 students per batch.}}
	\begin{CIErubrics}
		\Crubric{I}{Approval of the selected topic, formulation of Problem Statement and 	Objectives along with Synopsis submission}{10}
		\Crubric{II}{Demonstrate the skill and knowledge by applying appropriate tools/techniques to design solution specific to the problem.}{30}
		\Crubric{III}{Demonstrates the work carried out through experimental results, analysis
			and testing; Exhibits writing and communication skills through presentations
			and report writing.}{60}
	\end{CIErubrics}
	\begin{SEE@Theory@rubrics}
		\SrubricB{1}{Write Up}{20}
		\SrubricB{2}{Demonstration of Minor Project Work}{60}
		\SrubricB{3}{Viva voce}{20} 
	\end{SEE@Theory@rubrics}
\end{credit}
%-------------Internship----------------
\begin{credit}{Proj}
	\variantchar{Int}
	\Lcredit{0}
	\Pcredit{5}
	\CIEduration{2} 
	\SEEduration{3} 
	\CIELmarks{100}
	\SEELmarks{100}
	\commoninstCIE{The evaluation committee shall consist of Guide, Professor, Associate Professor/Assistant Professor. The	committee shall assess and evaluate the presentation and the progress reports.}
	\commoninstSEE{The SEE examination shall be conducted by an external examiner (domain expert) and an internal
		examiner. \textbf{Evaluation shall be carried out in batches, not exceeding 6 students per batch.}}
	\begin{CIErubrics}
		\Crubric{I}{Ability to comprehend the functioning/operating procedures of the
			Organization/Departments. Application of Engineering knowledge, Critical thinking and analysis to solve problems.}{40} 
		\Crubric{II}{Demonstrates skill competencies, Resource Management and Sustainability. Exhibits writing and communication skills through presentations
			and report writing.}{60}
	\end{CIErubrics}
	\begin{SEE@Theory@rubrics}
		\SrubricB{1}{Write Up}{20}
		\SrubricB{2}{Demonstration of Minor Project Work}{60}
		\SrubricB{3}{Viva voce}{20} 
	\end{SEE@Theory@rubrics}
\end{credit}


%-------------Major Project----------------
\begin{credit}{Proj}
	\variantchar{Maj}
	\Lcredit{0}
	\Pcredit{5}
	\CIEduration{2} 
	\SEEduration{3} 
	\CIELmarks{100}
	\SEELmarks{100}
	\commoninstCIE{The evaluation committee shall consist of Guide, Professor, Associate Professor/Assistant
		Professor. The committee shall assess and evaluate the presentation and the progress reports.}
	\commoninstSEE{Major Project SEE evaluation shall be conducted in two stages. This is initiated after fulfillment of submission of Project Report and CIE marks.\newline
		\textbf{Stage-1 Report Evaluation:} Evaluation of Project Report shall be carried out by Guide and an External 	examiner.\newline
		\textbf{Stage-2 Project Viva-voce:} Major Project Viva-voce examination is conducted after receipt of evaluation reports from Guide and External examiner.}
	\begin{CIErubrics}
		\Crubric{I}{Approval of the selected topic, formulation of Problem Statement and 	Objectives along with Synopsis submission}{10}
		\Crubric{II}{Demonstrate the skill and knowledge by applying appropriate tools/techniques to design solution specific to the problem.}{30}
		\Crubric{III}{Demonstrates the work carried out through experimental results, analysis and testing; Exhibits writing and communication skills through presentations and report writing.}{60}
	\end{CIErubrics}
	\begin{SEE@Theory@rubrics}
		\SrubricB{1}{Report evaluation: by Internal and external examiors for 100 marks each. \textbf{Final marks reduced to 50}.}{50}
		\SrubricB{2}{Viva-Voce: Jointly evaluated by Internal Guide \& External Examiners for 100 marks \textbf{Final marks reduced to 50}.}{50}
	\end{SEE@Theory@rubrics}
\end{credit}

%-----------Lecture Only Courses----------------
% 300
	\begin{credit}{LO}
	%\variantchar{a}
	\Lcredit{3}
	\CIEduration{2} 
	\SEEduration{3} 
	\CIELmarks{100}
	\SEELmarks{100}
	\begin{CIErubrics}
		\Crubric{1}{\textbf {QUIZZES:} Quiz will be conducted in online/offline mode. One QUIZ will be conducted \& will be evaluated for 10 marks. \textbf{This will be considered as final quiz marks}.}{10}
		\Crubric{2}{\textbf {TESTS:} Students will be evaluated in test through descriptive questions with different complexity levels (Revised Bloom’s Taxonomy Levels: Remembering, Understanding, Applying, Analyzing, Evaluating, and Creating). THREE tests will be conducted (Two regular tests and One optional Improvement test). Each test will be evaluated for 50 Marks, adding upto 100 Marks. \textbf{Final test marks will be reduced to 40}.}{40}
		\Crubric{3}{\textbf {EXPERIENTIAL LEARNING:} Students shall be evaluated based on the creativity, innovation, and practical implementation of their proposed solutions. The evaluation shall be conducted in two phases.  The implementation modalities for experiential learning activity shall be decided considering the nature of the course and its interdisciplinary requirements. \textbf{Each phase shall carry 25 marks, with the total evaluation carrying 50 marks.}}{50}
	\end{CIErubrics}
	\begin{SEE@Theory@rubrics}
		\SrubricB{1\space\&\space 2}{Unit - I: Question 1 or 2}{20}
		\SrubricB{3\space\&\space 4}{Unit - II: Question 3 or 4}{20}
		\SrubricB{5\space\&\space 6}{Unit - III: Question 5 or 6}{20}
		\SrubricB{7\space\&\space 8}{Unit - IV: Question 7 or 8}{20}
		\SrubricB{9\space\&\space 10}{Unit - V: Question 9 or 10}{20}
	\end{SEE@Theory@rubrics}
\end{credit}	

%200

\begin{credit}{LO}
	\variantchar{b}
	\Lcredit{2}
	\CIEduration{2} 
	\SEEduration{3} 
	\CIELmarks{50}
	\SEELmarks{50}
	\begin{CIErubrics}
		\Crubric{1}{\textbf {QUIZZES:} One Quiz will be conducted in online/offline mode \& will be evaluated for 10 Marks. \textbf{This will be considered as the final quiz marks.}}{10}
		\Crubric{2}{\textbf {TESTS:} Students will be evaluated in test, descriptive questions with different
			complexity levels (Revised Bloom’s Taxonomy Levels: Remembering, Understanding, Applying, Analyzing, Evaluating, and Creating). TWO tests will be conducted. Each test will be
			evaluated for 25 Marks, adding up to 50 Marks. \textbf{Final test marks will be reduced to 20 marks.}}{20}
		\Crubric{3}{\textbf {EXPERIENTIAL LEARNING:} \textbf {EXPERIENTIAL LEARNING:} Students shall be evaluated based on the creativity, innovation, and practical implementation of their proposed solutions. The evaluation shall be conducted in two phases.  The implementation modalities for experiential learning activity shall be decided considering the nature of the course and its interdisciplinary requirements. \textbf{Each phase shall carry 10 marks, with the total evaluation carrying 20 marks.}}{20}
	\end{CIErubrics}
	\begin{SEE@Theory@rubrics} 
	\SrubricA{1}{Objective type questions covering entire syllabus}{10}
	\SrubricB{1}{Unit - I: Compulsory}{12}
	\SrubricB{2\space\&\space 3}{Unit - II: Question 2 or 3}{14}
	\SrubricB{4\space\&\space 5}{Unit - III: Question 4 or 5}{14} 
	\end{SEE@Theory@rubrics}
\end{credit}	

%100
\begin{credit}{LO}
	\variantchar{c}
	\Lcredit{1}
	\CIEduration{2} 
	\SEEduration{2} 
	\CIELmarks{50}
	\SEELmarks{50}
	\begin{CIErubrics}
		\Crubric{1}{\textbf {QUIZZES:} One Quiz will be conducted in online/offline mode \& will be evaluated for 10 Marks. \textbf{This will be considered as the final quiz marks.}}{10}
		\Crubric{2}{\textbf {TESTS:} Students will be evaluated in test, descriptive questions with different complexity levels (Revised Bloom’s Taxonomy Levels: Remembering, Understanding, Applying, Analyzing, Evaluating, and Creating). TWO tests will be conducted. Each test will be evaluated for 50 Marks, adding upto 100 Marks. \textbf{Final test marks will be reduced to 20 marks.}}{20}
		\Crubric{3}{\textbf {EXPERIENTIAL LEARNING:} Students shall be evaluated based on the creativity, innovation, and practical implementation of their proposed solutions. The evaluation shall be conducted in two phases.  The implementation modalities for experiential learning activity shall be decided considering the nature of the course and its interdisciplinary requirements. \textbf{Each phase shall carry 20 marks, with the total evaluation carrying 40 marks. The final marks is reduced to 20 marks.}}{20}
	\end{CIErubrics}
	\begin{SEE@Theory@rubrics}
		\SrubricA{1}{Objective type questions (MCQs) covering the entire syllabus}{50}
	\end{SEE@Theory@rubrics}
\end{credit}	



%-----------Lecture with Tutorial Courses----------------
\begin{credit}{LT}
	%			\variantchar{1}
	\Lcredit{3}
	\Tcredit{1}
	\CIEduration{2} 
	\SEEduration{3} 
	\CIELmarks{100}
	\SEELmarks{100}
	%		\CIEPmarks{100}
	%		\SEEPmarks{100}
	\begin{CIErubrics}
		\Crubric{1}{\textbf{QUIZZES:} Quiz will be conducted in online/offline mode. ONE QUIZ will be conducted \& will be evaluated for 10 marks. \textbf{This will be considered as final quiz marks.}
		}{10}
		\Crubric{2}{\textbf{TESTS:} Students will be evaluated in test through descriptive questions with different complexity levels (Revised Bloom’s Taxonomy Levels: Remembering, Understanding, Applying, Analyzing, Evaluating, and Creating). THREE tests will be conducted (Two regular tests and One optional Improvement test). Each test will be evaluated for 50 Marks, adding upto 100 Marks. \textbf{Final test marks will be reduced to 40}.}{40}
		\Crubric{3}{\textbf {EXPERIENTIAL LEARNING:} Students shall be evaluated based on the creativity, innovation, and practical implementation of their proposed solutions. The evaluation shall be conducted in two phases.  The implementation modalities for experiential learning activity shall be decided considering the nature of the course and its interdisciplinary requirements. \textbf{Each phase shall carry 25 marks, with the total evaluation carrying 50 marks.}}{50}
	\end{CIErubrics}
	\begin{SEE@Theory@rubrics}[THEORY]
		\SrubricB{1\space\&\space 2}{Unit - I: Question 1 or 2}{20}
		\SrubricB{3\space\&\space 4}{Unit - II: Question 3 or 4}{20}
		\SrubricB{5\space\&\space 6}{Unit - III: Question 5 or 6}{20}
		\SrubricB{7\space\&\space 8}{Unit - IV: Question 7 or 8}{20}
		\SrubricB{9\space\&\space 10}{Unit - V: Question 9 or 10}{20}
	\end{SEE@Theory@rubrics}
\end{credit}



%-----------Lecture with Practicle Courses----------------
	\begin{credit}{LP}
	%\variantchar{a}
	\Lcredit{3}
	\Pcredit{1}
	\CIEduration{2+2} 
	\SEEduration{3+3} 
	\CIELmarks{100}
	\CIEPmarks{50}
	\SEELmarks{100}
	\SEEPmarks{50}
	\begin{CIErubrics}
		\Crubric{1}{\textbf {QUIZZES:} Quiz will be conducted in online/offline mode. ONE quiz will be conducted \& will be evaluated for 10 marks. \textbf{This will be considered as final quiz marks.}
		}{10}
		\Crubric{2}{\textbf{TESTS:} Students will be evaluated in test through descriptive questions with different complexity levels (Revised Bloom’s Taxonomy Levels: Remembering, Understanding, Applying, Analyzing, Evaluating, and Creating). THREE tests will be conducted (Two regular tests and One optional Improvement test). Each test will be evaluated for 50 Marks, adding upto 100 Marks. \textbf{Final test marks will be reduced to 40}.}{40}
		\Crubric{3}{\textbf {EXPERIENTIAL LEARNING:} Students shall be evaluated based on the creativity, innovation, and practical implementation of their proposed solutions. The evaluation shall be conducted in two phases.  The implementation modalities for experiential learning activity shall be decided considering the nature of the course and its interdisciplinary requirements. \textbf{Each phase shall carry 25 marks, with the total evaluation carrying 50 marks.}}{50}
		\Crubric{4}{\textbf {LAB:} Conduction of laboratory exercises, lab report, observation, and analysis (30 Marks), PBL progress review (10 marks) and lab test (10 Marks) adding up to 50 Marks. \textbf{The final lab marks is for 50.}}{50}
	\end{CIErubrics}
	\begin{SEE@Theory@rubrics}
		\SrubricB{1\space\&\space 2}{Unit - I: Question 1 or 2}{20}
		\SrubricB{3\space\&\space 4}{Unit - II: Question 3 or 4}{20}
		\SrubricB{5\space\&\space 6}{Unit - III: Question 5 or 6}{20}
		\SrubricB{7\space\&\space 8}{Unit - IV: Question 7 or 8}{20}
		\SrubricB{9\space\&\space 10}{Unit - V: Question 9 or 10}{20}
	\end{SEE@Theory@rubrics}
	\begin{SEE@Practical@rubrics}
		\Srubric{1}{Evaluation will be conducted through problem-/project-based learning activities specific to the program or laboratory. The assessment rubrics will be mapped to the respective course outcomes and shared with students along with the problem statements.}{50}
		%	\Srubric{2}{Conduction of the Experiments}{20}
		%	\Srubric{3}{Viva}{20}
	\end{SEE@Practical@rubrics}
\end{credit}
	

%-----------Practicle Only Courses----------------
\begin{credit}{PO}
	%\variantchar{a}
	\Pcredit{1}
	\CIEduration{2} 
	\SEEduration{2} 
	\CIEPmarks{50}
	\SEEPmarks{50}
	%	\CIELmarks{50}
	%	\SEELmarks{50}
	\begin{CIErubrics}
		\Crubric{1}{\textbf{QUIZZES:} Quizzes will be conducted in online/offline mode. TWO quizzes will be conducted \& Each Quiz will be evaluated for 10 Marks. \textbf{The average of two quizzes will be the final quiz marks.}}{10}
		\Crubric{2}{\textbf{TESTS:} Students will be evaluated in test, descriptive questions with different complexity levels (Revised Bloom’s Taxonomy Levels: Remembering, Understanding, Applying, Analyzing, Evaluating, and Creating). TWO tests will be conducted. Each test will be evaluated for 50 Marks, adding upto 100 Marks. \textbf{Final test marks will be reduced to 20 marks.}}{20}
		\Crubric{3}{\textbf{EXPERIENTIAL LEARNING:} Students will be evaluated for their creativity and practical implementation of the problem. Case study-based teaching learning (10), Program specific requirements (10), Video based seminar/presentation/demonstration (20) adding upto 40 marks. \textbf{The final EL marks is reduced to 20 marks.}}{20}
	\end{CIErubrics}
	\begin{SEE@Practical@rubrics}
		\Srubric{1}{Objective type questions (MCQs) covering the entire syllabus}{50}
		%	\Srubric{2}{Unit - I: Compulsory}{16}
		%	\Srubric{3\space\&\space 4}{Unit - II: Compulsory}{16}
	\end{SEE@Practical@rubrics}
\end{credit}
\begin{credit}{PO}
	\variantchar{b}
	\Pcredit{1}
	\CIEduration{2} 
	\SEEduration{2} 
	\CIEPmarks{50}
	\SEEPmarks{50}
	%	\CIELmarks{50}
	%	\SEELmarks{50}
	\begin{CIErubrics}
		\Crubric{1}{\textbf{Modules Assessment:} Modules will be assessed for 30 marks}{30}
		\Crubric{2}{\textbf{TESTS:} Two tests will be conducted for 10 marks each adding to 20 marks}{20}
	\end{CIErubrics}
	\begin{SEE@Practical@rubrics}
		\Srubric{1}{Objective type questions (MCQs) covering the entire syllabus}{50}
		%	\Srubric{2}{Unit - I: Compulsory}{16}
		%	\Srubric{3\space\&\space 4}{Unit - II: Compulsory}{16}
	\end{SEE@Practical@rubrics}
\end{credit}


\end{definecredits}

